\chapter{Wilson Fermions (WF)}

In Section~8, it was shown that the lattice action is not unique.  One has the
freedom to add an arbitrary number of irrelevant operators to the action, as
these do not change the continuum limit.  Wilson's solution to the doubling
problem was to add a dimension five operator $a r \bar\psi \square \psi$,
whereby the extra fifteen species at $p_\mu = \pi$ get a mass proportional to
$r/a$ [56].  The Wilson action is
\begin{align}
A^W &= m_q \sum_x \bar\psi(x)\psi(x) \nonumber\\
&\quad + \frac{1}{2a}\sum_{x,\mu} \bar\psi(x)\gamma_\mu \bigl[ U_\mu(x)\psi(x+\hat{\mu}) - U_\mu^\dagger(x-\hat{\mu})\psi(x-\hat{\mu}) \bigr] \nonumber\\
&\quad - \frac{r}{2a}\sum_{x,\mu} \bar\psi(x) \bigl[ U_\mu(x)\psi(x+\hat{\mu}) - 2\psi(x) + U_\mu^\dagger(x-\hat{\mu})\psi(x-\hat{\mu}) \bigr] \nonumber\\
&= \frac{(m_q a + 4r)}{a}\sum_x \bar\psi(x)\psi(x) \nonumber\\
&\quad + \frac{1}{2a}\sum_x \bar\psi \bigl[ (\gamma_\mu - r) U_\mu(x)\psi(x+\hat{\mu}) - (\gamma_\mu + r) U_\mu^\dagger(x-\hat{\mu})\psi(x-\hat{\mu}) \bigr] \nonumber\\
&\equiv \sum_{x,y} \bar\psi^L_x M^W_{xy} \psi^L_y
\end{align}
where the interaction matrix $M^W$ is usually written as
\begin{equation}
M^W_{x,y}[U]\,a = \delta_{xy} - \kappa \sum_\mu \bigl[ (r - \gamma_\mu) U_{x,\mu}\delta_{x,y-\mu} + (r + \gamma_\mu) U^\dagger_{x-\mu,\mu}\delta_{x,y+\mu} \bigr]
\end{equation}
with the rescaling
\begin{align}
\kappa &= 1/(2m_q a + 8r) \nonumber\\
\psi^L &= \sqrt{m_q a + 4r}\ \psi = \psi/\sqrt{2\kappa} \,.
\end{align}

In this form the Wilson action has the following interpretation as a
Statistical Mechanics model.  The local term acts like an inertia keeping the
fermion at the same site while the non-local term makes the fermion hop to the
nearest neighbor site with strength $\kappa$.  During this hop the fermion
acquires a twist by $(\gamma_\mu - r)$ in spin space and by $U_\mu$ in color
space.  The case $r = 1$ is special as the spin twist $1 \pm \gamma_\mu$ is a
projection operator of rank 2,
\begin{align}
\Bigl( \frac{1 \pm \gamma_\mu}{2} \Bigr)^2 &= \frac{1 \pm \gamma_\mu}{2} \nonumber\\
\mathrm{Tr} \Bigl( \frac{1 \pm \gamma_\mu}{2} \Bigr) &= 2 \,.
\end{align}
Later we shall see that for $r = 1$ the dispersion relation has a single
branch, \textit{i.e.} there are no doublers in the free field limit.

From Eq.~10.3, one notes that for the free theory the quark mass is given in
terms of the lattice parameters $\kappa$ and $r$ as
\begin{equation}
m_q a = \frac{1}{2\kappa} - 4r \equiv \frac{1}{2\kappa} - \frac{1}{2\kappa_c}
\end{equation}
with a zero at $\kappa = \kappa_c \equiv 1/8r$.  For the interacting theory
($U_\mu(x) \neq 1$) we will continue to define $m_q a = 1/2\kappa - 1/2\kappa_c$
with the proviso that $\kappa_c$ depends on $a$.  The renormalization of
$\kappa_c$ independent of $\kappa$ implies that the quark mass has both
multiplicative and additive renormalizations.  This is a consequence of the
explicit breaking of chiral symmetry by the irrelevant term proportional to
$r$.  Wilson's fix for doublers comes with a high price tag --- a hard breaking
of chiral symmetry at $O(a)$.

The free propagator in momentum space for Wilson fermions is
\begin{equation}
S_F(p) = M_W^{-1}(p) = \frac{a}{\,1 - 2\kappa \sum_\mu \bigl( r \cos p_\mu a - i\gamma_\mu \sin p_\mu a \bigr)}
\end{equation}
Due to the dimension 5 operator (terms proportional to $r$) the 15 extra states
at $p_\mu a = \pi$ get masses of order $2rn/a$, where $n$ is the number of
components of $p_\mu a = \pi$.

The operator $M$ satisfies the following relations:
\begin{align}
\gamma_5 M_W^\dagger \gamma_5 &= M_W, \nonumber\\
\gamma_5 S_F^\dagger (x,y) \gamma_5 &= S_F (y,x) \ , \nonumber\\
M_W^\dagger (\kappa, r) &= M_W (-\kappa, -r) \ .
\end{align}

The first two state that the ``hermiticity'' property of $M^N$ is preserved.
The second equation relates a quark propagator from $x \to y$ to the antiquark
propagator from $y \to x$.  This important identity, called hermiticity or
$\gamma_5$ invariance, leads to a significant simplification in numerical
calculations as discussed in Section~17.1.  The adjoint in $S_F^\dagger$ is
with respect to the spin and color indices at each site.  The last relation
shows that $M^W$ is not anti-hermitian due to the Wilson $r$ term.

An analysis similar to Eq.~7.9 for the naive fermion action, shows that the
pole mass derived from Eq.~10.6 is different from the bare mass and given by
\begin{equation}
m_q^{\mathrm{pole}} a = r(\cosh Ea - 1) + \sinh Ea \,.
\end{equation}
This shows, as expected, that the discretization corrections in spectral
quantities occur at $O(a)$.

\section{Properties of Wilson Fermions}

I give a brief summary of the important properties of Wilson fermions.

\begin{itemize}
\item The doublers are given a heavy mass, $2 r n /a$, and decouple in the
continuum limit.

\item Chiral symmetry is broken explicitly.  The derivation of axial Ward
identities, using the invariance under the transformation given in Eq.~7.11,
have the generic form
\begin{equation}
\Bigl\langle \frac{\partial \delta\mathcal{S}}{\partial \theta}\,\mathcal{O} \Bigr\rangle = \Bigl\langle \frac{\partial \delta\mathcal{O}}{\partial \theta} \Bigr\rangle \,.
\end{equation}
For WF, the variation of the action under an axial transformation is
\begin{equation}
\frac{\partial \delta\mathcal{S}}{\partial \theta} = \partial_\mu A_\mu - 2 m P + r a X,
\end{equation}
where $X$ is an additional term coming from the variation of the Wilson $r$
term [57].  Thus, in general, all relations based on axial WI will have
corrections at order $ra$ and involve mixing with operators that would normally
be absent due to the chiral symmetry.

\item Calculation of the matrix element $\langle 0| T[A_\mu V_\nu V_\rho] |0\rangle$
shows that the ABJ anomaly is correctly reproduced in the continuum
limit [58].  Karsten and Smit show, by an explicit 1-loop calculation, that
while each of the extra 15 states contribute terms proportional to $r$, the
total contribution from all sixteen states is independent of $r$ and equals the
ABJ anomaly.

\item The quark mass gets both additive and multiplicative renormalization.

\item The zero of the quark mass is set by $\kappa_c$ in Eq.~10.5.  There are
two ways to calculate $\kappa_c$ at any given $a$.  (i) Assume the chiral
relation $M_\pi^2 \propto m_q$, calculate the pion mass as a function of
$1/2\kappa$, and extrapolate it to zero.  The value of $\kappa$ at which the
pion becomes massless is, by definition, $\kappa_c$.  (ii) Calculate the quark
mass through the ratio $\langle \partial_\mu A_\mu(x) P(0)\rangle / \langle
P(x) P(0)\rangle$ (based on the axial Ward identity) as a function of
$1/2\kappa$ and extrapolate to zero.  The two estimates can differ by
corrections of $O(a)$.  Both these calculations requires a statistical average
over gauge configurations and an extrapolation in $1/2\kappa$.  Consequently,
on an individual configuration the zero mode of the matrix $M$ will not occur
at $\kappa_c$ but will have a distribution about $\kappa_c$.

In the full theory $\det M$ suppresses configurations with zero modes,
however, this protection is lost in QQCD.  Nevertheless, the problem of zero
modes at $\kappa < \kappa_c$ on certain configurations, called
``exceptional'' configurations, needs to be addressed in both theories.

Near these zero-modes the quark propagator becomes singular even though one is
in the physical region $\kappa < \kappa_c$.  Consequently, in a large enough
statistical sample one will hit ``exceptional'' configurations, and hadronic
correlations functions on these will show fluctuations by orders of magnitude
(depending on small the quark mass is) from the rest.  Since this pathology is
a feature of Wilson-like actions and not of the gauge configurations
themselves, one cannot simply remove these configurations from the statistical
sample.  The problem of such zero-modes increases with decreasing quark mass
and $\beta$.  Two practical approaches used so far have been to (i) simulate at
$m_q > m_q^{*}$ where $m_q^{*}$ is sufficiently large such that the spread in
the position of the zero modes does not give rise to large fluctuations in the
correlation function, and (ii) for fixed $m_q$ increase the lattice volume as
it decreases the probability of hitting exceptional configurations.  These
fixes do not help if one is interested in simulating QCD near the chiral limit,
\textit{i.e.}, for physical light quarks.  For this new ideas are needed.  A
more detailed analysis of this problem and a possible solution is given
in [59,60].

\item The spin and flavor degrees of freedom are in one-to-one correspondence
with continuum Dirac fermions.  Thus, the construction of interpolating field
operators is straightforward, for example, $\bar\psi \gamma_5 \psi$ and $\bar
\psi \gamma_i \psi$ are interpolating operators for pseudoscalar and vector
mesons just as in the continuum (see Section~17).

\item The Wilson term changes the discretization errors to $O(a)$.

\item \textbf{Feynman Rules} and details of gauge fixing for the plaquette
gauge action and Wilson fermions are given in [36].

\item \textbf{Symmetries of $S_F$}: The transformation of the propagator under
the discrete symmetries is the same as for naive fermions.  For a given
background configuration $U$ these are
\begin{align}
\mathcal{P}: \qquad S_F(x,y,[U]) &\to \gamma_4\,S_F(x^P,y^P,[U^P])\,\gamma_4 \nonumber\\
\mathcal{T}: \qquad S_F(x,y,[U]) &\to \gamma_4\gamma_5\,S_F(x^T,y^T,[U^T])\,\gamma_5\gamma_4 \nonumber\\
\mathcal{C}: \qquad S_F(x,y,[U]) &\to \gamma_4\gamma_2\,S_F^{\mathrm{T}}(x,y,[U^C])\,\gamma_2\gamma_4 \nonumber\\
\mathcal{H}: \qquad S_F(x,y,[U]) &\to \gamma_5\,S_F^\dagger(y,x,[U])\,\gamma_5 \nonumber\\
\mathcal{CPH}: \quad S_F(x,y,[U]) &\to \mathcal{C}\gamma_4\,S_F^\dagger(y^P,x^P,[U^C])\,\gamma_4\mathcal{C}^{-1}
\end{align}
These relations are very useful in determining the properties of correlation
functions (real versus imaginary, even or odd in $p$, \textit{etc.}) as
discussed in Section~17.

\item The dispersion relation, Eq.~10.8, has the following two solutions
\begin{equation}
e^{Ea} = \frac{(ma + r) \pm \sqrt{1 + 2mar + m^2 a^2}}{1+r}
\end{equation}
For $r = 0$ and $m$ small (naive fermion action), the two solutions are
\begin{align}
E &= m + O(m^2 a^2) \nonumber\\
E &= -m + \frac{i\pi}{a} + O(m^2 a^2) \,,
\end{align}
where the second solution corresponds to the pole in the Euclidean propagator
at $p_4 = \pi$ and is the temporal doubler.  For $r = 1$ this doubler is
removed, as expected, and the single root $Ea = \log(1 + ma)$ is associated
with the physical particle.  For a general $r \neq 1$ the second root is given
by
\begin{equation}
E = -m + \frac{1}{a} \Bigl( i\pi + \log \frac{1-r}{1+r} \Bigr) + \ldots
\end{equation}
Such branches which are not associated with the physical state or with doublers
are called ghost branch.  I use a very restrictive definition of a doubler ---
a solution which has $|E| = 0$ for $m = \vec{p} = 0$.  Again from Eq.~10.14 it
is clear that this ghost branch is pushed to infinity for $r = 1$.  In general,
if the action couples $(n+1)$ time slices, then the dispersion relation has
$n$ branches.  The roots can be real, complex-conjugate pairs, and/or have an
imaginary part $i\pi/a$ as illustrated above.  Such additional branches can be
a nuisance for improved actions, for example the $D234$ action discussed in
Section~13.2.  The goal is to obtain a doubler and ghost free theory.  This can
be done by a combination of tuning the parameters (for example choosing $r = 1$
in the Wilson case) and/or going to anisotropic lattices $a_t \ll a_s$ as
advocated by Lepage and collaborators [107].

\item \textbf{Conserved Vector Current}: The Wilson action is invariant under
the global transformation defined in Eq.~7.10.  To derive the associated
conserved current (for degenerate masses) we use the standard trick of
calculating the variation of the action under a space-time dependent
$\theta(x)$ and associating the coefficient of $\theta(x)$ in $\delta S$ with
$\partial_\mu J_\mu$.  The change in the action under such an infinitesimal
transformation is
\begin{align*}
\delta S &= \kappa \sum_{x,\mu} \bar\psi(x)(\gamma_\mu - r) U_\mu(x) \psi(x + \hat{\mu}) \exp\bigl( i\theta(x) - i\theta(x + \mu) \bigr) \\
&\quad - \kappa \sum_{x,\mu} \bar\psi(x + \mu)(\gamma_\mu + r) U_\mu^\dagger (x) \psi(x) \exp\bigl( i\theta(x + \mu) - i\theta(x) \bigr)
\end{align*}
which, to first order in $\theta$, is
\[
-\sum_{x,\mu} \Bigl[ \bar\psi(x)(\gamma_\mu - r) U_\mu(x) \psi(x + \hat{\mu}) + \bar\psi(x + \mu)(\gamma_\mu + r) U_\mu^\dagger (x) \psi(x) \Bigr] \Bigl[ i \frac{\partial \theta}{\partial x_\mu} \Bigr] .
\]
The conserved current, obtained after integration by parts, is
\begin{equation}
V^{c}_\mu = \bar\psi(x) (\gamma_\mu - r) U_\mu (x) \psi (x + \hat{\mu}) + \bar\psi(x + \mu)(\gamma_\mu + r) U_\mu^\dagger (x) \psi (x) \,.
\end{equation}
$V^{c}_\mu$ is hermitian and reduces to the symmetrized version of the 1-link
vector current for $r = 0$.  In many applications like decay constants one uses
the local (flavor) currents defined by
\begin{align}
2V_\mu(x) &= \bar\psi(x) \gamma_\mu \frac{\lambda^a}{2} U_\mu (x) \psi (x + \hat{\mu}) + \mathrm{h.c.} \nonumber\\
2A_\mu(x) &= \bar\psi(x) \gamma_5\gamma_\mu \frac{\lambda^a}{2} U_\mu (x) \psi (x + \hat{\mu}) + \mathrm{h.c.} \,,
\end{align}
which are not conserved, and consequently have associated non-trivial
renormalization factors $Z_V$ and $Z_A$ which have to included when calculating
matrix elements.

\item \textbf{The chiral condensate}: As explained in [57], the chiral
condensate $\langle \bar\psi \psi \rangle$ is not simply related to the trace
of the Wilson quark propagator $\langle S_F(0,0) \rangle$.  The breaking of
chiral symmetry by the $r$ term introduces contact terms that need to be
subtracted non-perturbatively from $S_F(0,0)$.  This has not proven practical.
The methods of choice are to use either the continuum Ward Identity
\begin{equation}
\langle \bar\psi\psi \rangle^{\mathrm{WI}} \equiv \langle 0| S_F(0,0) |0\rangle = \lim_{m_q \to 0} m_q \int d^4x\, \langle 0| P(x)P(0) |0\rangle,
\end{equation}
or the lattice version of the Gell-Mann, Oakes, Renner relation [61]
\begin{equation}
\langle \bar\psi\psi \rangle^{\mathrm{GMOR}} = \lim_{m_q \to 0} - \frac{f_\pi^2 M_\pi^2}{4 m_q} \,.
\end{equation}
A comparison of the efficacy of the two methods is discussed in [62].

\item \textbf{Operator Mixing}: There is, in general, a mixing of operators
with different chirality in the Callen Symanzik equations.  This is true even
at $\kappa_c$ and is a consequence of the explicit breaking of chiral symmetry.
This mixing of operators poses serious problems in the the calculation of
matrix elements of the weak interactions Hamiltonian, \textit{i.e.}, the
lattice analyses of Ward identities and obtaining amplitudes with the correct
chiral behavior is more involved.  For example the chiral behavior of matrix
elements of the tree-level 4-fermion operators (relevant to the extraction of
$B_K$, $B_7$, and $B_8$) is not the same as in the continuum theory.  The
reason being that the Ward Identities associated with the spontaneously broken
$SU_A(n_f)$ receive corrections at $O(a)$ in the lattice theory.  This issue
will be discussed in more detail by Prof.~Martinelli.
\end{itemize}

In Section~13 I will discuss how the $O(a)$ errors in $WF$ can, with very
little extra effort, be reduced to $O(a^2)$ or higher.  For the time being let
me use this drawback of Wilson fermions to motivate the formulation of
staggered fermions.  This formulation has $O(a^2)$ errors, retains enough
chiral symmetry to give the correct chiral behavior for amplitudes involving
Goldstone pions, but at the expense of a 4-fold increase in the number of
flavors.
